% =========================================================
% Paper text: asymptotic corollary for Theorem 3
% =========================================================
% Dr Yang Azzollini — 21 April 2026
%
% Intended placement: immediately after the statement of
% Theorem 3 in the paper (wherever two-step efficiency is
% introduced). This text complements the finite-sample
% statement already formalised in Lean.
%
% Style: matches the existing paper tone — spare, technical,
% no overselling. Uses the same notation conventions as
% Yangs_theorems_summary.pdf.
% =========================================================


\begin{corollary}[Asymptotic attainment of the two-step bound]
\label{cor:twostep-asymptotic}
Let $\{\hat\kappa_N\}$ be a sequence of first-step estimates of $\kappa = \rho^2$
obtained by substituting the Corollary~1.1 estimator into
$\hat\kappa = (\hat\sigma_{12}^{(1)})^2 / (\hat\sigma_1^2 \hat\sigma_2^2)$,
indexed by sample size $N$. Suppose $\hat\kappa_N \xrightarrow{\PP} \kappa_0$
for some $\kappa_0 \in [0,1]$ as $N \to \infty$. Then the two-step estimator
$\hat\sigma_{12}^{(2)}$ based on $\hat\kappa_N$ attains the variance bound of
Theorem~3 at $\kappa = \kappa_0$:
\[
  \Var\!\left(\hat\sigma_{12}^{(2)} \,\big|\, \tau, \hat\kappa_N\right)
  \;\xrightarrow{\PP}\;
  \min_{w : T(w) = 1}\; \Var\!\left(\hat\sigma_{12}(w) \,\big|\, \tau, \kappa_0\right)
  \qquad \text{as } N \to \infty.
\]
\end{corollary}

\begin{proof}[Proof sketch]
The KKT system of Theorem~2 defines the optimal weights $w^*(\kappa)$ as the
unique solution of a $(K{+}1) \times (K{+}1)$ linear system whose coefficients
are continuous in $\kappa$. Since the system matrix is non-singular on
$[0,1]$ (as established in the proof of Theorem~2), the solution
$w^*(\kappa)$ is continuous in $\kappa$ by Cramer's rule, and the
variance $V(w^*(\kappa))$ is continuous in $\kappa$ as a quadratic form.
Continuous mapping applied to $\hat\kappa_N \xrightarrow{\PP} \kappa_0$
gives the result.
\end{proof}

\begin{remark}[Scope of the Lean~4 verification]
The finite-sample fixed-$\kappa$ form of Theorem~3 is machine-verified in
Lean~4 (see \texttt{ZHY\_Opt.two\_step\_efficiency}, commit
\texttt{116f0df}). Corollary~\ref{cor:twostep-asymptotic} adds a standard
continuity-plus-consistency argument whose probabilistic content —
convergence in probability of $\hat\kappa_N$, continuous mapping — lies
outside the deterministic-$\tau$ setting of the Lean formalisation, and
is established on paper by the standard arguments above.
\end{remark}
